\documentclass[a4paper,12pt]{article}
\usepackage[fontset=none]{ctex}
\usepackage{geometry}
\usepackage{graphicx}
\usepackage{amsmath}
\usepackage{enumitem}
\usepackage{setspace}
\usepackage{fontspec}

% 字体设置（与 docx 模板一致）
\setCJKmainfont{Songti SC}          % 宋体 - 正文默认
\setCJKsansfont{Heiti SC}            % 黑体 - 一级标题
\setCJKmonofont[AutoFakeBold]{STFangsong}          % 仿宋 - 注意事项
\setmainfont{Times New Roman}        % 英文默认

% 页面设置
\geometry{left=2.5cm, right=2.5cm, top=2.5cm, bottom=2.5cm}

% 去掉页码
\pagestyle{empty}

\begin{document}

% ========== 封面 ==========
\begin{center}
    \vspace*{2cm}
    {\fontsize{36pt}{43pt}\selectfont \textbf{物理实验预习报告}}
    \vspace{3cm}
\end{center}

\begin{center}
    \fontsize{16pt}{24pt}\selectfont
    \textbf{实验名称：密立根油滴实验\quad} \\[0.5em]
    \textbf{指导教师：\_\_\_\_\_\_\_\_\_\_\_\_\_\_\_\_\_\_\_\_\_\_\_\_\_}
\end{center}

\vspace{2cm}

\begin{center}
    \fontsize{14pt}{28pt}\selectfont
    \textbf{周次：\_\_\_\_\_\_\_\_\_\_\_\_\_\_\_\_\_\_\_\_\_\_\_\_\_\_\_\_\_\_\_\_\_} \\
    \textbf{姓名：\_\_\_\_\_\_\_\_\_\_\_\_\_\_\_\_\_\_\_\_\_\_\_\_\_\_\_\_\_\_\_\_\_} \\
    \textbf{学号：\_\_\_\_\_\_\_\_\_\_\_\_\_\_\_\_\_\_\_\_\_\_\_\_\_\_\_\_\_\_\_\_\_} \\[0.5em]
    \textbf{实验日期: \_\_\_\_年\_\_\_\_月\_\_\_\_日\quad 星期\_\_\_上/下午}
\end{center}

\vspace{1.5cm}

\begin{center}
    \fontsize{14pt}{20pt}\selectfont
    浙江大学物理实验教学中心
\end{center}

\newpage

% ========== 正文 ==========

{\fontsize{14pt}{20pt}\selectfont \textbf{1. 实验综述}}

\vspace{0.3cm}

{\fontsize{12pt}{18pt}\selectfont

密立根油滴实验是近代物理学的经典实验之一，首次精确测定了基本电荷$e$，证实了电荷的量子性——所有带电体的电量均是最小基本电荷$e$的整数倍。

实验原理：带电油滴在两平行极板间同时受重力和静电力作用。实验采用\textbf{静态平衡法}与\textbf{匀速下降法}相结合：（1）\textbf{静态平衡}：调节极板电压$U$使油滴悬停，此时静电力与重力平衡，即$qU/d = mg$；（2）\textbf{匀速下降}：撤去电压后油滴在重力作用下加速，最终达到终端速度$v_g$匀速下落，由斯托克斯定律（引入康宁汉修正）得到油滴修正半径$r'$，进而求出质量$m$；（3）由平衡电压$U$、板间距$d$和质量$m$计算每颗油滴的带电量$q = mgd/U$；（4）对多颗油滴的电荷量用\textbf{公约数法}求最大公约数，即基本电荷量$e$。

实验装置为密立根油滴仪：平行极板（板间距$d$约5\,mm）、显微镜、CCD成像系统、显示屏、喷雾器及计时器。实验选取带电量小于10$e$、平衡电压约200\,V、下落时间在10\,s/mm以内的清晰油滴。

}

\vspace{0.5cm}

{\fontsize{14pt}{20pt}\selectfont \textbf{2. 实验重点}}

\vspace{0.3cm}

{\fontsize{12pt}{18pt}\selectfont

掌握静态平衡法与匀速下降法相结合的测量原理；理解康宁汉修正公式对微小油滴半径的修正意义；学会由终端速度$v_g$和平衡电压$U$计算电荷量$q$，并熟练运用公约数法提取基本电荷量$e$，验证电荷量子化。

}

\vspace{0.5cm}

{\fontsize{14pt}{20pt}\selectfont \textbf{3. 实验难点}}

\vspace{0.3cm}

{\fontsize{12pt}{18pt}\selectfont

在CCD显示屏上准确甄别并跟踪合适油滴（带电量适中，不受布朗运动显著干扰）；精确判断油滴达到匀速下落状态以减小计时误差；对各油滴电荷量进行整数化时正确估计所带基本电荷数$n$，避免倍数判断的失误；对测量数据合理估算不确定度。

}

\vfill

% ========== 注意事项 ==========
\noindent\rule{\textwidth}{0.4pt}

{\CJKfontspec[AutoFakeBold]{STFangsong}\fontsize{12pt}{18pt}\selectfont
\textbf{注意事项：}
\begin{enumerate}[label=\arabic*., leftmargin=2em, nosep]
    \item 用PDF格式上传"预习报告"，文件名：学生姓名+学号+实验名称+周次。
    \item "预习报告"必须递交在"学在浙大"的本课程的对应实验项目的"作业"模块内。
    \item "预习报告"还须拷贝到"实验报告"中（便以教师批改）。
\end{enumerate}
}

\vspace{0.5cm}
\begin{center}
    {\CJKfontspec[AutoFakeBold]{STFangsong}\fontsize{12pt}{18pt}\selectfont \textbf{浙江大学物理实验教学中心制}}
\end{center}

\end{document}
